\subsection{Quotient spaces}

Let X be a path-connected topological space equipped with a proper (right) action of a discrete group G. Set Y = X/G and denote by \(f: X \to Y\) the canonical map. If \(g \in G\) and c: I → X is a path in X, we denote by \(g^{*}c\) the path \(t \mapsto c(t) \cdot g\) and \(g^{*}[c]\) its strict homotopy class.

Let \(o\) be a point of X. For every \(g \in G\) , let \(\beta_{g}\) be the class of a path joining \(o \cdot g\) to o in X. For every \(g \in G\) , let \(X^{g}\) be the set of points \(x \in X\) such that \(x \cdot g = x\) ; for every path-connected component j of \(X^{g}\) , let \(a_{j}\) be a point of j and let \(\gamma_{j}\) be the class of a path in X joining \(a_{j}\) to o. Let \(\nu: F(G) \to \pi_{1}(Y, f(o))\) be the unique homomorphism of groups such that \(\nu(g) = f_{*}(\beta_{g})\) for \(g \in G\) . Let N: \(\pi_{1}(X, o) * F(G) \to \pi_{1}(Y, f(o))\) be the unique group homomorphism which coincides with \(\pi_{1}(f, o)\) on \(\pi_{1}(X, o)\) and with \(\nu\) on \(F(G)\) .

PROPOSITION 3. --- Suppose that X is delaceable. The homomorphism N is then surjective and its kernel is the smallest distinguished subgroup of \(\pi_{1}(\mathrm{X},o)*\mathrm{F}(\mathrm{G})\) containing the elements

\[
r _ {2} (k, v) = [ k ] ^ {- 1} v [ k ] (\beta_ {k} ^ {- 1} k ^ {*} (v) ^ {- 1} \beta_ {k})\tag{\((R_{2})\)}
\]

\[
\text {for } k \in \mathrm{G} \text { and } v \in \pi_ {1} (\mathrm{X}, o);\tag{\((R_{3}^{\prime})\)}
\]

\[
r _ {3} ^ {\prime} (k, j) = [ k ] (\beta_ {k} ^ {- 1} k ^ {*} (\gamma_ {j}) ^ {- 1} \gamma_ {j})
\]

for \(k\in \mathrm{G}\) and \(j\in \pi_0(\mathrm{X}^k)\)

\[
r _ {3} ^ {\prime \prime} (k, h) = [ k h ] ^ {- 1} [ k ] [ h ] (\beta_ {h} ^ {- 1} h ^ {*} (\beta_ {k} ^ {- 1}) \beta_ {k h})\tag{\((\mathrm{R}_{3}^{\prime \prime})\)}
\]

\[
\text {for } k \text { and } h \in G.
\]

The groupoid morphism \(\varpi(f)\) induces a morphism

\[
\varpi^ {\prime \prime} (f) \colon \varpi (\mathrm{X}) / \mathrm{G} \to \varpi (\mathrm{Y})
\]

which is an isomorphism by theorem 3 (IV, p. 403), since the space X is assumed to be delaceable. The proposition then follows from II, p. 211, prop. 6.

The three corollaries below follow from the immediately corresponding corollaries of Prop. 6 of II, p. 211.

COROLLARY 1. --- Suppose that X is delaceable and that the group G is generated by the stabilizers of the points of X. The canonical morphism \(\pi_{1}(f,o)\colon\pi_{1}(\mathrm{X},o)\to\pi_{1}(\mathrm{Y},f(o))\) is then surjective. In particular, if X is simply connected by arcs, then so is Y.

Remark. --- If X is delaceable, then so is Y (IV, p. 349, prop. 8). Since a connected delaceable space is simply connected by arcs if and only if it is simply connected (IV, p. 344, corollary 1 of theorem 1), we thus recover prop. 11 of I, p. 137.

Example 1. --- Let X be a path-connected, delaceable, and separated topological space, and let a be a point of X. Let n be an integer \(\geqslant 2\) and let Y be the quotient of the space \(X^{n}\) by the action of the group \(S_{n}\) acting by permutation of the factors; denote by \(f: X^{n} \to Y\) the canonical map; denote \(g: X \to Y\) the map \(x \mapsto f(x, a, \ldots, a)\) . It follows from the proposition that, for every i, the homomorphism \(\pi_{1}(g, a)\) from \(\pi_{1}(X, a)\) into \(\pi_{1}(Y, g(a))\) is surjective and that its kernel is the derived subgroup of \(\pi_{1}(X, a)\) . In particular, the group \(\pi_{1}(Y, g(a))\) is abelian.

COROLLARY 2. --- Suppose that X is delaceable and that the group G acts freely on X. There exists a unique group homomorphism \(p \colon \pi_1(\mathrm{Y}, f(o)) \to \mathrm{G}\) whose kernel contains the image of \(\pi_1(\mathrm{X}, o)\) and such that \(p(\mathsf{N}(g)) = g\) for all \(g \in \mathrm{G}\) . Moreover, \(\pi_1(\mathrm{X}, o) \to \pi_1(\mathrm{Y}, f(o)) \xrightarrow{p} \mathrm{G}\) is an extension of G by \(\pi_1(\mathrm{X}, o)\) .

COROLLARY 3. --- Suppose that X is simply connected by arcs. The map from G to \(\pi_1(Y, f(o))\) which associates to \(g \in G\) the path class \(f_*(\beta_g)\) is a surjective group homomorphism; its kernel is the subgroup of G generated by the stabilizers of the points of X.

